% Representación analítica correlacionada de alfa: acción, convergencia e
% intervalos.
\subsubsection{Representación analítica de la coordenada de alfa}
\label{subsec:alpha-lector-completo-rev08}

La prolongación analítica se especifica mediante todos sus coeficientes.
Sea \(\mathsf s\in\mathcal X_\alpha^{\mathrm{enr}}\). Antes
de la normalización decimal, el mismo estado ya tiene la precoordenada
\begin{equation}
 a(\mathsf s):=p(\mathsf s)+e_0(\mathsf s)-f(\mathsf s)
                 -\kappa_{\mathrm{per}}(\mathsf s)\in I_\alpha,
 \label{eq:alpha-precoordenada-comun}
\end{equation}
con \(p=\pi_{\mathrm{HMT}}-3\),
\(e_0=e_{\mathrm{HMT}}-2\) y
\(f=\varphi_{\mathrm{HMT}}-1\). Equivalentemente,
\(a=\pi_{\mathrm{HMT}}+e_{\mathrm{HMT}}-
\varphi_{\mathrm{HMT}}-4-\kappa_{\mathrm{per}}\). No se introduce aquí
\(\alpha\) convencional: \(p,e_0,f\) son los tres caracteres regionales ya
generados y \(\kappa_{\mathrm{per}}\) es la lectura racional del registro
terminal. La clausura entera aplica a
\eqref{eq:alpha-precoordenada-comun} la normalización de
\eqref{eq:alpha-recurrencia}; la coordenatización que sigue expresa sobre la misma
precoordenada una representación analítica compatible. Su grafo de
dependencia recibe \((p,e_0,f,\kappa_{\mathrm{per}})\) directamente: no
contiene una arista \(\alpha_A\to\lambda\), aunque la igualdad demostrada más
abajo permita identificar después ambos nombres de representación.

Escribamos \(q=1/729\), sea \(F_{\mathsf s}\) el cuerpo ordenado generado por
las coordenadas de \(\mathsf s\), y pongamos
\(E_{9,\mathsf s}(X)=E_{21/22}^{(9)}(\mathsf s;X)\). Definimos el coeficiente
de enlace
\begin{equation}
 \lambda(\mathsf s):=
 -\frac{E_{9,\mathsf s}(a(\mathsf s))
         \bigl(1-q\,a(\mathsf s)^3\bigr)}{a(\mathsf s)^{10}}.
 \label{eq:alpha-lambda-estado}
\end{equation}
Todos sus argumentos pertenecen al estado antes del reconocimiento
metrológico. En particular, \(\lambda\) no se obtiene ajustando una cadena
decimal objetivo. La ecuación~\eqref{eq:alpha-lambda-estado} debe leerse con
su tipo exacto: una vez que la primera representación del estado ha producido
\(a(\mathsf s)\), fija la única prolongación de la clase geométrica que se
define a continuación. No constituye una segunda selección independiente de
\(a(\mathsf s)\), sino una representación analítica correlacionada del mismo
estado enriquecido.

En efecto, para \(\mu\in F_{\mathsf s}\) pongamos
\[
 E_{\mathsf s,\mu}(X)
 :=E_{9,\mathsf s}(X)+\mu\frac{X^{10}}{1-qX^3}.
\]
Como \(0<a(\mathsf s)<10^{-2}\) y \(qa(\mathsf s)^3<1\), se tiene
\begin{equation}
 E_{\mathsf s,\mu}(a(\mathsf s))=0
 \quad\Longleftrightarrow\quad
 \mu=\lambda(\mathsf s).
 \label{eq:alpha-unicidad-lambda-clase-geometrica}
\end{equation}
Así, la semilla de coeficientes no contiene libertad una vez fijados el estado
y la clase de prolongaciones \(q\)-geométricas. Esta unicidad es relativa a la
clase declarada; no afirma que el jet de orden nueve seleccione por sí solo una
cola entre todas las series analíticas posibles.

La acción local sobre bloques de tres coeficientes es
\begin{equation}
 \mathcal C_\alpha:F_{\mathsf s}^{3}\longrightarrow F_{\mathsf s}^{3},
 \qquad \mathcal C_\alpha(u,v,w)=q(u,v,w),
 \qquad c^{(0)}=(\lambda(\mathsf s),0,0).
 \label{eq:alpha-accion-coeficientes}
\end{equation}
Por tanto, para todo \(n\geq0\),
\begin{equation}
 b_{10+3n}=q^n\lambda(\mathsf s),\qquad
 b_{11+3n}=b_{12+3n}=0.
 \label{eq:alpha-recurrencia-coeficientes}
\end{equation}
Ésta es la regla que calcula cada coeficiente posterior al grado nueve. No
queda un parámetro libre en cada prolongación.

Para \(M\geq0\), sea
\begin{equation}
 E_{\mathsf s,M}(X):=E_{9,\mathsf s}(X)
 +\lambda(\mathsf s)\sum_{n=0}^{M}q^nX^{10+3n}.
 \label{eq:alpha-truncamientos-completos}
\end{equation}
La función completa es
\begin{equation}
 E_{\mathsf s,\infty}(X)
 :=E_{9,\mathsf s}(X)
 +\lambda(\mathsf s)\frac{X^{10}}{1-X^3/729}.
 \label{eq:alpha-funcion-completa-explicita}
\end{equation}

Para formular el sistema proyectivo sin tipos implícitos, fijamos
\(d_M=12+3M\) y definimos la fibración de jets
\begin{equation}
 \mathfrak J_M:=
 \coprod_{\mathsf s\in\mathcal X_\alpha^{\mathrm{enr}}}
 \{\mathsf s\}\times F_{\mathsf s}[X]/(X^{d_M+1}).
 \label{eq:alpha-espacio-jets-rev08}
\end{equation}
Si \(M'\ge M\), la restricción es
\begin{equation}
 \tau_{M',M}(\mathsf s,[P]_{d_{M'}})
 :=(\mathsf s,[P]_{d_M}).
 \label{eq:alpha-truncamiento-jets-rev08}
\end{equation}

\begin{proposition}[Naturalidad de la recurrencia analítica]
\label{prop:alpha-naturalidad-recurrencia-explicita}
Sea
\[
 \Gamma_M(\mathsf s)
 =\bigl(\mathsf s,[E_{\mathsf s,M}]_{d_M}\bigr).
\]
Si \(M'\geq M\), el truncamiento \(\tau_{M',M}\) satisface
\begin{equation}
 \tau_{M',M}\circ\Gamma_{M'}=\Gamma_M,
 \qquad
 \Gamma_{E21}(\mathsf s)=\varprojlim_M\Gamma_M(\mathsf s).
 \label{eq:alpha-naturalidad-explicita}
\end{equation}
La proyección sobre los tres grados nuevos es precisamente
\(\mathcal C_\alpha^{M+1}(c^{(0)})\).
\end{proposition}

\begin{proof}
La ecuación~\eqref{eq:alpha-recurrencia-coeficientes} fija el bloque de
grados \(10+3n,11+3n,12+3n\). Al pasar de \(M\) a \(M+1\) sólo se añade
el bloque siguiente; truncarlo recupera literalmente el polinomio anterior.
La compatibilidad para \(M'\geq M\) se obtiene por composición. Así, el
límite inverso no es una sección escogida del jet desnudo, sino la órbita
única de \(c^{(0)}\) bajo \(\mathcal C_\alpha\).
\end{proof}

\subsubsection{Convergencia uniforme, raíz simple e intervalos racionales}
\label{subsec:alpha-cotas-completas-rev08}

Los cilindros regionales y el registro periódico proporcionan cotas mediante
aritmética de intervalos racionales. Para cada coordenada
\[
 c\in\{\pi_{\mathrm{HMT}},e_{\mathrm{HMT}},\varphi_{\mathrm{HMT}}\},
\]
sea \(P_c\) su prefijo nonádico certificado de mil cifras. El dato utilizado
no es el racional \(P_c\) como si fuese el valor completo, sino el cilindro
semiabierto y su cierre exterior
\[
 \mathcal C_c^{(1000)}=[P_c,P_c+\varepsilon),
 \qquad
 \overline{\mathcal C}_c^{(1000)}=[P_c,P_c+\varepsilon],
 \qquad \varepsilon=10^{-1000}.
\]
La aritmética dirigida opera conservadoramente con esos cierres. Si
\(P_\pi,P_e,P_\varphi\) denotan los tres prefijos, se obtiene
\begin{equation}
\begin{aligned}
 a_-&:=P_\pi+P_e-(P_\varphi+\varepsilon)-4-
       \kappa_{\mathrm{per}},\\
 a_+&:=(P_\pi+\varepsilon)+(P_e+\varepsilon)-P_\varphi-4-
       \kappa_{\mathrm{per}},\\
 A_{1000}&:=[a_-,a_+],
 \qquad a(\mathsf s)\in[a_-,a_+]=A_{1000},
 \qquad \operatorname{diam}A_{1000}=3\varepsilon.
\end{aligned}
\label{eq:alpha-propagacion-colas-rev08}
\end{equation}
La extensión intervalar natural de
\eqref{eq:alpha-lambda-estado} transporta después \(A_{1000}\), el cilindro
de \(\pi_{\mathrm{HMT}}\) y la cota racional dirigida de \(\delta\) a un
intervalo cerrado \(\Lambda_{1000}\ni\lambda(\mathsf s)\); no se congela
ninguna de las tres colas. En particular,
\begin{equation}
\begin{aligned}
 \frac{7297352569283800997285105472380662}{10^{36}}
 &<a(\mathsf s)<
 \frac{7297352569283800997285105472380663}{10^{36}},\\
 -8.711\mathbin{\times}10^{-6}
 &<\lambda(\mathsf s)<-8.709\mathbin{\times}10^{-6}.
\end{aligned}
 \label{eq:alpha-cotas-precoordenada-lambda}
\end{equation}
Estas desigualdades usan prefijos certificados con su cota racional de
cola; no reciben una tabla de alfa. El verificador del paquete las reproduce
desde las tres salidas nonádicas y desde la coordenatización terminal.

Sea
\[
 r:=\frac{10^{-6}}{729}<1.372\mathbin{\times}10^{-9}.
\]
Para \(0\leq x\leq10^{-2}\), la cola posterior a \(M\) satisface
\begin{equation}
 \left|E_{\mathsf s,\infty}(x)-E_{\mathsf s,M}(x)\right|
 \leq
 8.711\mathbin{\times}10^{-26}\,\frac{r^{M+1}}{1-r}.
 \label{eq:alpha-cota-cola-uniforme}
\end{equation}
Además,
\begin{equation}
 \left|\frac{\mathrm d}{\mathrm dx}
  \left(\lambda\frac{x^{10}}{1-x^3/729}\right)\right|
 \leq 8.711\mathbin{\times}10^{-6}
 \left(\frac{10\,10^{-18}}{1-r}
 +\frac{3\,10^{-24}/729}{(1-r)^2}\right)
 <8.712\mathbin{\times}10^{-23}.
 \label{eq:alpha-cota-derivada-cola}
\end{equation}
La cola de las derivadas posterior al bloque \(M\) admite, además, la
estimación dependiente de la profundidad
\begin{equation}
 \sup_{0\le x\le10^{-2}}
 \left|E'_{\mathsf s,\infty}(x)-E'_{\mathsf s,M}(x)\right|
 \le 8.711\mathbin{\times}10^{-24}\,r^{M+1}
 \left(\frac{13+3M}{1-r}+\frac{3r}{(1-r)^2}\right)
 \xrightarrow[M\to\infty]{}0.
 \label{eq:alpha-cota-resto-derivadas-rev08}
\end{equation}

\begin{theorem}[Función límite y raíz simple única de la coordenatización correlacionada]
\label{thm:alpha-raiz-completa-explicita}
La sucesión \(E_{\mathsf s,M}\) converge uniformemente, junto con sus
derivadas, a \(E_{\mathsf s,\infty}\) en
\(\overline I_\alpha=[0,10^{-2}]\). Se cumple
\begin{equation}
 E_{\mathsf s,\infty}(a(\mathsf s))=0,
 \qquad
 \inf_{x\in\overline I_\alpha}
 E'_{\mathsf s,\infty}(x)>6.239.
 \label{eq:alpha-raiz-y-derivada-completas}
\end{equation}
Por consiguiente, \(a(\mathsf s)\) es la raíz simple única de la función
completa en \(I_\alpha\).
\end{theorem}

\begin{proof}
La razón de dos términos consecutivos de la cola está acotada por \(r<1\),
lo que da \eqref{eq:alpha-cota-cola-uniforme}; la derivación término a
término y la suma de las dos series geométricas dan
\eqref{eq:alpha-cota-derivada-cola} y
\eqref{eq:alpha-cota-resto-derivadas-rev08}. La prueba del jet establece
\(E'_{9,\mathsf s}>6.24\) en \(\overline I_\alpha\); por tanto, la
perturbación completa deja la derivada por encima de \(6.239\).
Finalmente, al sustituir \eqref{eq:alpha-lambda-estado} en
\eqref{eq:alpha-funcion-completa-explicita}, los dos sumandos evaluados en
\(a(\mathsf s)\) se cancelan exactamente. La función es estrictamente
creciente; su única raíz es, pues, simple y coincide con esa precoordenada.
\end{proof}

Para hacer efectiva la identificación a toda profundidad, las evaluaciones
dirigidas de los extremos certifican primero
\[
 E_{\mathsf s,\infty}(0)<0<
 E_{\mathsf s,\infty}(10^{-2}),
 \qquad
 E'_{\mathsf s,\infty}\geq\mu:=6239/1000>c:=6
 \quad\hbox{en }[0,10^{-2}].
\]
Por continuidad existe una raíz en el intervalo y, por la segunda desigualdad,
es única. Sólo después de fijar esta existencia, la monotonía y la invariante
de que la horquilla contiene la raíz se aplica el rescate por el teorema del
valor medio. Partimos de \(D_0=[0,10^{-2}]\). Si
\(D_j=[r_j,s_j]\), sea \(m_j=(r_j+s_j)/2\), y sea
\([F_j^-,F_j^+]\) una inclusión racional dirigida de
\(E_{\mathsf s,\infty}(m_j)\), obtenida propagando los cilindros de
\(\pi_{\mathrm{HMT}},e_{\mathrm{HMT}},\varphi_{\mathrm{HMT}}\),
\(\delta\) y \(\lambda\), con
\begin{equation}
 0\leq F_j^+-F_j^-\leq \frac{c}{4}\,(s_j-r_j).
 \label{eq:alpha-anchura-evaluacion-rescate-rev08}
\end{equation}
Cuando \(F_j^+<0\) se conserva la mitad derecha;
cuando \(F_j^->0\), la mitad izquierda. Si
\(F_j^-\leq0\leq F_j^+\), no se intenta decidir si el punto medio es
exactamente la raíz. Sea \(w_j=s_j-r_j\). La cota uniforme
\(E'_{\mathsf s,\infty}\geq c=6\) y el teorema del valor medio implican que
cualquier raíz contenida en \(D_j\) pertenece a
\begin{equation}
 D_{j+1}=D_j\cap
 \left[m_j-\frac{w_j}{4},
             m_j+\frac{w_j}{4}\right].
 \label{eq:alpha-rescate-derivada-rev08}
\end{equation}
En efecto, sea \(\xi_j\) la raíz y sea
\(y_j=E_{\mathsf s,\infty}(m_j)\in[F_j^-,F_j^+]\). Como el recinto contiene
cero y satisface \eqref{eq:alpha-anchura-evaluacion-rescate-rev08}, se tiene
\(|y_j|\leq cw_j/4\). El teorema del valor medio proporciona
\(|m_j-\xi_j|\leq |y_j|/c\leq w_j/4\), lo que prueba
\eqref{eq:alpha-rescate-derivada-rev08}. El argumento incluye el caso
\(y_j=0\): la raíz está entonces en el punto medio, pero el algoritmo no
necesita reconocer esa igualdad para conservarla.

Si el recinto disponible no satisface
\eqref{eq:alpha-anchura-evaluacion-rescate-rev08}, se refinan primero los
cilindros fuente y se repite la evaluación; no se declara un signo ni se
acepta una contracción insuficiente. En cualquiera de las tres ramas el nuevo
intervalo conserva la raíz y tiene diámetro a lo sumo \(w_j/2\). Para cada
profundidad objetivo finita, el refinamiento dirigido permite satisfacer la
cota de anchura sin comparar de forma no certificada un número real con cero.
La monotonía demuestra inductivamente
\begin{equation}
 a(\mathsf s)\in D_{j+1}\subseteq D_j,\qquad
 \operatorname{diam}D_{j+1}\leq
 \tfrac12\operatorname{diam}D_j,
 \qquad \operatorname{diam}D_j\longrightarrow0,
 \label{eq:alpha-biseccion-racional}
\end{equation}
donde la desigualdad es una cota de contracción y no una igualdad impuesta.
El certificado ejecutable aplica
esta regla a veinticuatro niveles en base mil y comprueba que el cierre
exterior completo \(A_{1000}\), no sólo su extremo inferior, queda dentro de
cada cilindro expresado. Sus pruebas negativas rechazan un recinto demasiado
ancho, un recinto desordenado, un punto que no sea el medio, un dominio sin
cambio de signo y una horquilla degenerada. La prolongación coinductiva permite
reemplazar mil por cualquier profundidad requerida.

Concretamente, para \(S_N=1000^N\) el verificador calcula
\begin{equation}
 j_N^-:=\lfloor S_Na_-\rfloor,
 \qquad j_N^+:=\lfloor S_Na_+\rfloor
 \label{eq:alpha-indices-cilindro-dirigidos-rev08}
\end{equation}
y sólo certifica el nivel si \(j_N^-=j_N^+\). Esta igualdad, que incluye el
extremo superior del cierre exterior, prueba
\(A_{1000}\subset[j_N^-/S_N,(j_N^-+1)/S_N)\) y evita atribuir a la celda
izquierda un extremo situado exactamente sobre su frontera derecha.

Sea \(C_N^A(\mathsf s)\) el cilindro entero de
\eqref{eq:alpha-cilindros-via-a-exactos}. Elegimos recursivamente
\(m(N)>m(N-1)\) de modo que
\(\operatorname{diam}D_{m(N)}<1000^{-N}/4\), y definimos
\begin{equation}
I_{B,N}:=D_{m(N)}\cap C_N^A(\mathsf s).
 \label{eq:alpha-intervalos-b-completos}
\end{equation}
Definimos
\(\ell_N:=\inf I_{B,N}\) y \(u_N:=\sup I_{B,N}\). Ambos son racionales;
el intervalo es no vacío porque los dos factores contienen
\(a(\mathsf s)\). Puesto que el cilindro es semiabierto, \(u_N\) puede no
pertenecer a \(I_{B,N}\). La evaluación en ese extremo se entiende mediante
la extensión continua de \(E_{\mathsf s,\infty}\) al cierre exterior. Así,
\begin{equation}
 I_{B,N+1}\subseteq I_{B,N}\subseteq C_N^A(\mathsf s),
 \qquad \operatorname{diam}I_{B,N}\longrightarrow0,
 \qquad E_{\mathsf s,\infty}(\ell_N)\leq0\leq
 E_{\mathsf s,\infty}(u_N).
 \label{eq:alpha-inclusion-signos-completa}
\end{equation}
La cota~\eqref{eq:alpha-cota-cola-uniforme} permite escoger un
truncamiento que conserve el signo en cada extremo donde la función completa
no se anula. Si un extremo coincide con la raíz, se utiliza la identidad de
la función completa y la pertenencia ya demostrada; no se atribuye ese mismo
signo a sus truncamientos. En efecto, el resto exacto es
\[
 E_{\mathsf s,\infty}(x)-E_{\mathsf s,M}(x)
 =\lambda(\mathsf s)x^{10}
   \frac{(qx^3)^{M+1}}{1-qx^3}.
\]
En el dominio positivo considerado, donde \(\lambda(\mathsf s)<0\), se
obtiene en la raíz
\[
 E_{\mathsf s,M}(a(\mathsf s))
 =-\lambda(\mathsf s)a(\mathsf s)^{10}
   \frac{(qa(\mathsf s)^3)^{M+1}}{1-qa(\mathsf s)^3}>0
 \qquad(M<\infty).
\]
Por ello, la certificación de los signos débiles en
\eqref{eq:alpha-inclusion-signos-completa} se refiere a
\(E_{\mathsf s,\infty}\). Su evaluación utiliza la función completa o el
truncamiento acompañado del resto con signo. Esto cubre tanto el ínfimo
perteneciente como el supremo eventualmente excluido del cilindro.

\begin{proposition}[Cuadrados de compatibilidad y cilindros comunes]
\label{prop:alpha-cuadrados-compatibilidad}
Las restricciones de palabras y de representaciones analíticas conmutan:
\[
\begin{array}{ccc}
 \mathcal X_\alpha^{\mathrm{enr}}&
 \xrightarrow{\ \mathsf T_{\alpha,N'}\ }&\mathcal W_{N'}\\[1mm]
 \Big\Vert&&\Big\downarrow\rho_{N',N}\\[1mm]
 \mathcal X_\alpha^{\mathrm{enr}}&
 \xrightarrow{\ \mathsf T_{\alpha,N}\ }&\mathcal W_N
\end{array}
\qquad
\begin{array}{ccc}
 \mathcal X_\alpha^{\mathrm{enr}}&
 \xrightarrow{\ \Gamma_{M'}\ }&\mathfrak J_{M'}\\[1mm]
 \Big\Vert&&\Big\downarrow\tau_{M',M}\\[1mm]
 \mathcal X_\alpha^{\mathrm{enr}}&
 \xrightarrow{\ \Gamma_M\ }&\mathfrak J_M .
\end{array}
\]
La colección \(I_{B,N}\) proporciona la conexión efectiva entre ambos
cuadrados y satisface materialmente
\(I_{B,N}\subseteq C_N^A\) para todo \(N\).
\end{proposition}

\begin{proof}
El primer cuadrado es la compatibilidad de la normalización entera; el segundo
es la proposición~\ref{prop:alpha-naturalidad-recurrencia-explicita}. La
inclusión se sigue de la definición~\eqref{eq:alpha-intervalos-b-completos},
y la no vacuidad y los signos se siguen del
teorema~\ref{thm:alpha-raiz-completa-explicita}. Ningún paso identifica la
raíz de un truncamiento con la de una cola arbitraria.
\end{proof}

Definimos ahora, sin dejar implícito el lector de prefijos,
\begin{equation}
 \alpha_B(\mathsf s):=
 \operatorname*{arg\,zero}_{x\in I_\alpha}E_{\mathsf s,\infty}(x),
 \qquad
 \operatorname{pref}_N(x):=
 \frac{\lfloor1000^Nx\rfloor}{1000^N},
 \qquad
 \alpha_{B,N}:=\operatorname{pref}_N(\alpha_B).
 \label{eq:alpha-definicion-b-y-prefijos-rev08}
\end{equation}

\begin{theorem}[Igualdad proyectiva de las dos representaciones correlacionadas]
\label{thm:alpha-dos-publicaciones-completas}
Para todo \(N\), la vía entera y la vía analítica dan el mismo prefijo:
\begin{equation}
 \boxed{\alpha_{A,N}=\operatorname{pref}_N(\alpha_A)
 =\operatorname{pref}_N(\alpha_B)=\alpha_{B,N}.}
 \label{eq:alpha-dos-vias-cada-nivel}
\end{equation}
En el límite,
\begin{equation}
 \boxed{\alpha_A
 =\operatorname*{arg\,zero}_{x\in I_\alpha}
       E_{\mathsf s,\infty}(x)
 =\alpha_B=\alpha_{\mathrm{HMT}}.}
 \label{eq:alpha-dos-vias-limite-coinductivo}
\end{equation}
\end{theorem}

\begin{proof}
La normalización entera representa exactamente \(a(\mathsf s)\). El
teorema~\ref{thm:alpha-raiz-completa-explicita} demuestra que la coordenatización
analítica representa ese mismo punto como raíz simple única. Los
cilindros semiabiertos \(C_N^A\) fijan un prefijo canónico y
\(I_{B,N}\subseteq C_N^A\) contiene \(\alpha_B=\alpha_A\). Por tanto, sus
prefijos coinciden para cada \(N\); el anidamiento y los diámetros tendentes a
cero dan la igualdad de los límites. El reconocimiento
metrológico ocurre después.
\end{proof}

La identidad demostrada enlaza dos representaciones del mismo estado
mediante la recurrencia de coeficientes y los cilindros compatibles.
La función analítica es una representación correlacionada de la
precoordenada ya generada, no una segunda selección causalmente
independiente. El jet, la clase geométrica y las coordenadas del estado
conservan sus funciones respectivas; el reconocimiento convencional es posterior.
